\section{Use of AI in the Research and Formalization}
\label{app:ai-assistance}

This appendix is a historical account of our workflow in 2025--2026.
As of October 2026, we consider it plausible that an agent using GPT-6 Astra,
given this manuscript with its Rocq excerpts (including those in
Appendix~\ref{app:formal-interface-excerpts}) removed and access to Rocq and
the stated dependencies, could independently recreate the full formalization
from scratch for at most US\$200 in model usage.
We have not performed this experiment; both its success and its cost remain
conjectural.

OpenAI Codex was used extensively in both the Rocq development and the
preparation of this manuscript.
Its contributions included implementations, proof scripts, auxiliary lemma
statements, mathematical arguments and counterexamples, explanations of
library interfaces, and drafts and revisions of the exposition.
The authors supplied the research objectives, developed the initial
specifications, directed and reviewed the subsequent work, and retain
responsibility for the definitions, assumptions, results, and text.
The account below draws on the authors' recollection, the Git histories of the
artifact and manuscript, and retained development conversations.
These records support a chronology and concrete examples of collaboration;
they do not support a precise division of credit by line count.

\subsection{Human design and early formalization}

The cryptographic objective preceded the intensive AI-assisted proof work.
In 2025, the authors formalized the approximate homomorphic-encryption
interface, correctness properties, and IND-CPA and IND-CPAD security games
largely by hand.
The artifact already contained a statement of the noise-flooding security
reduction in August 2025, although its proof and some quantitative parameters
were still placeholders.
The intended security result therefore guided the later proof development;
the final interfaces, assumptions, and numerical bound were refined during
that development.

The relational logic similarly began with human-written judgments and proposed
rules in LaTeX, followed by manual encodings of the judgments as Rocq
propositions over SSProve semantics.
These early definitions and rule statements were starting points for
verification, and several required revision before soundness could be proved.

\subsection{Joint development of logic soundness}

After a break, substantial AI-assisted work resumed in May 2026.
Much of the initial soundness work proceeded through short, interactive
exchanges: the authors selected a rule or probability lemma, Codex proposed a
proof or decomposed it into auxiliary obligations, and the authors questioned
the statements, discussed the argument, and redirected the next step.
Codex implemented proof scripts and used compiler feedback to revise them.
The collaboration covered the probability results underlying the logic as
well as the program-level rules, including conditional KL composition,
Pinsker's inequality, and discrete-Gaussian analysis.

Two changes illustrate how this process affected the mathematics.
First, in June, discussion of sequencing exposed a hidden-mixture problem:
omitting a shared random intermediate value from the transcript can invalidate
the proposed coordinatewise conditional-KL bounds.
Codex supplied a counterexample during the discussion, and the subsequent
rule design retained such intermediate values as explicit coordinates with
zero KL budget.
These zero entries preserve the conditioning information needed by later
coordinates without increasing the final numerical loss.
Keeping them explicit was part of repairing the sequence rule.

Second, the authors raised the behavior of the judgments on programs with
missing probability mass and asked whether the existing completion
construction could be used for the Pythagorean judgment.
Together with Codex, they developed and proved the completed sequencing rules
and the bridge to additive error.
The resulting definitions account for missing subdistribution mass with an
explicit \(\bot\) outcome, as described in \Cref{sec:pyth-toolkit,sec:artifact}.
Both episodes involved changes to proposed statements and definitions,
followed by proofs of the revised rules.

\subsection{Compiler construction and verification}

The authors identified the need to expose a bounded number of selected oracle
calls in an otherwise arbitrary adaptive adversary, and supplied the intended
algorithmic behavior and semantic preservation requirement.
The trace compiler was developed interactively in early June 2026 from this
design, building on a partial trace/replay implementation that the authors
supplied for review.
Codex filled in substantial parts of the implementation, helped formulate the
correctness theorem, and constructed its proof through intermediate lemmas.
The authors reviewed the construction and corrected proposed restrictions on
the theorem, including an unnecessary bounded-call hypothesis for
same-package semantic correctness: calls remaining after the selected prefix
are handled by the residual linked package.

The discussions also clarified SSProve's operation resolution, raw-code
linking, and package composition interfaces.
The migration to SSProve's nominal extension had already begun in late May,
before the trace compiler, with AI assistance in adapting the constructions
and repairing interfaces.
The quantitative compiler rule was a further joint development: its statement
was revised as the logic evolved, and Codex supplied substantial proofs for
the adaptive replacement theorem and its probability dependencies.
By late June, some of these proof tasks were delegated for extended execution
between author checkpoints, alongside continued interactive review.

\subsection{Application to noise flooding and game-hop exposition}

For the application, the authors supplied the security target, an existing
pen-and-paper proof outline, and the strategy of accumulating per-decryption
KL costs before a single conversion to statistical distance.
From late June through July, Codex carried out substantial implementation and
proof work on the concrete reduction, invariants, local decryption comparison,
and composition of exact and quantitative transitions.
This included extended tasks using Codex's persistent-goal workflow, with
progress records and intervening author discussion of assumptions and proof
strategy.

The detailed game presentations in \Cref{sec:trace-compiler} emerged from
this formal proof and several rounds of manuscript feedback.
Codex drafted the intermediate representations and their explanations; the
authors asked for a conventional cryptographic reduction, explicit endpoints,
and a clear distinction between the original game and its compiled form.
In particular, in July the authors explicitly proposed the seven-stage
\(H_0,\ldots,H_6\) organization, separating the original IND-CPAD experiment,
its linked representation, and its compiled representation before the paid
replacement step.
Later joint revisions distinguished this implementation-level factorization
from the three conceptual games in \Cref{eq:conceptual-game-chain}.
The final presentation therefore reflects both AI drafting and specific human
choices about the proof's structure and explanation.

\subsection{Review, checking, and manuscript preparation}

AI suggestions also introduced errors that required author intervention.
For example, during the application proof an AI-added assumption equating
encodings through different inverse charts was challenged by the authors
because it conflicted with the intended origin-centered charts.
The assumption was removed from the proof route.
The replacement argument compares Gaussian centers in a common chart and then
uses shared postprocessing; its specification and proof were developed through
further discussion and formal verification.
This episode illustrates the separate roles of checking a proof relative to
its assumptions and reviewing whether those assumptions express the intended
cryptographic setting.

Rocq kernel checking and inspection of \texttt{Print Assumptions} were central
to accepting the resulting proofs.
The final trusted base and remaining inherited assumption are documented in
\Cref{sec:tcb}; AI-generated proof scripts are subject to the same checks as
human-written scripts.
Human review concentrated especially on the scheme and game interfaces,
theorem statements, assumptions, and the interpretation of the checked result.
